%% Font size %%
\documentclass[11pt]{article}

%% Load the custom package
\usepackage{Mathdoc}

%% Numéro de séquence %% Titre de la séquence %%
\renewcommand{\centerhead}{Fractions II - Éval 1 : Diviseurs}

%% Spacing commands %%
\renewcommand{\baselinestretch}{1} \setlength{\parindent}{0pt}

\begin{document}

\phantom{0}

\entetedevoirs{20}

\begin{center}
\duree{1 heures} 
\total{50 points}
\coefficient{1}
\calculatrice{0}
\brouillon
\soin
\recherche 
\copieseparee{0}
\end{center}

\begin{exercicedevoir}[5][Décomposer en produits de facteurs premiers]
\textit{Exemple :} $36 = 6 \times 6 = 2 \times 3 \times 2 \times 3$
\begin{enumerate}
\item $24 = \ldots\ldots\ldots \times \ldots\ldots\ldots = \ldots\ldots\ldots \times \ldots\ldots\ldots \times \ldots\ldots\ldots \times \ldots\ldots\ldots $
\item $45 = \ldots\ldots\ldots \times \ldots\ldots\ldots = \ldots\ldots\ldots \times \ldots\ldots\ldots \times \ldots\ldots\ldots $
\item $54 = \ldots\ldots\ldots \times \ldots\ldots\ldots = \ldots\ldots\ldots \times \ldots\ldots\ldots \times \ldots\ldots\ldots \times \ldots\ldots\ldots$
\item $72 = \ldots\ldots\ldots \times \ldots\ldots\ldots = \ldots\ldots\ldots \times \ldots\ldots\ldots \times \ldots\ldots\ldots \times \ldots\ldots\ldots \times \ldots\ldots\ldots$
\item $63 = \ldots\ldots\ldots \times \ldots\ldots\ldots = \ldots\ldots\ldots \times \ldots\ldots\ldots \times \ldots\ldots\ldots $
\end{enumerate}
\end{exercicedevoir}


\begin{exercicedevoir}[5][Déterminer la liste des diviseurs I]
\begin{multicols}{2}
\begin{enumerate}[itemsep=1em]
\item Compléter le tableau suivant et faire la liste de tous les diviseurs de 66.$\medskip$\\$\begin{array}{|c|c|c|}
\hline
\text{Facteur n$^\circ$ 1} & \text{Facteur n$^\circ$ 2} & \text{Produit donnant } 66 \\
\hline
\ldots & \ldots& \ldots \\
\hline
2 & 33& \ldots \\
\hline
3 & 22& \ldots \\
\hline
\ldots & 11& 66 \\
\hline
\end{array}
$\\
\item Compléter le tableau suivant et faire la liste de tous les diviseurs de 48.$\medskip$\\$\begin{array}{|c|c|c|}
\hline
\text{Facteur n$^\circ$ 1} & \text{Facteur n$^\circ$ 2} & \text{Produit donnant } 48 \\
\hline
\ldots & 48& \ldots \\
\hline
\ldots & \ldots& 48 \\
\hline
3 & 16& 48 \\
\hline
\ldots & 12& 48 \\
\hline
\ldots & 8& 48 \\
\hline
\end{array}
$\\
\item Compléter le tableau suivant et faire la liste de tous les diviseurs de 96.$\medskip$\\$\begin{array}{|c|c|c|}
\hline
\text{Facteur n$^\circ$ 1} & \text{Facteur n$^\circ$ 2} & \text{Produit donnant } 96 \\
\hline
\ldots & \ldots& 96 \\
\hline
\ldots & 48& \ldots \\
\hline
3 & 32& \ldots \\
\hline
\ldots & 24& \ldots \\
\hline
6 & 16& \ldots \\
\hline
\ldots & \ldots& 96 \\
\hline
\end{array}
$\\
\item Compléter le tableau suivant et faire la liste de tous les diviseurs de 28.$\medskip$\\$\begin{array}{|c|c|c|}
\hline
\text{Facteur n$^\circ$ 1} & \text{Facteur n$^\circ$ 2} & \text{Produit donnant } 28 \\
\hline
\ldots & \ldots& \ldots \\
\hline
2 & 14& \ldots \\
\hline
4 & 7& \ldots \\
\hline
\end{array}
$\\
\end{enumerate}
\end{multicols}
\end{exercicedevoir}



\begin{exercicedevoir}[5][Déterminer la liste des diviseurs II]
\begin{enumerate}
\item Écrire la liste de tous les diviseurs de 99. \\ \\ \\ \\ 
\item Écrire la liste de tous les diviseurs de 88. \\ \\ \\ \\
\item Écrire la liste de tous les diviseurs de 84. \\ \\ \\ \\
\item Écrire la liste de tous les diviseurs de 30. \\ \\ \\ \\
\end{enumerate}
\end{exercicedevoir}

\begin{exercicedevoir}[5][Résoudre le problème suivant]
\textit{Toute trace de recherche (schémas, essais, calculs, tableaux, etc.) sera valorisée dans l'évaluation.}\\ \\
Un garçon de café doit répartir 36 croissants et 24 pains au chocolat dans des corbeilles. \\
Chaque corbeille doit avoir le même contenu. \\
Quelles sont les répartitions possibles ?
\end{exercicedevoir}


\end{document}


%%% Local Variables:
%%% mode: LaTeX
%%% TeX-master: t
%%% TeX-master: t
%%% End:

